% ============================================================
% 第1章 赛题理解与总体技术路线
% 【负责人：组长/全员讨论后由组长执笔】
% 内容要求：设计思路（赛题“6.1 提交内容”明确要求说明文档含设计思路）
% 已写好赛题理解的客观部分，技术路线叙述与流程图需全组补充。
% ============================================================
\section{赛题理解与总体技术路线}

\subsection{赛题任务理解}

任务一要求基于 G7 易流提供的 500 台车辆 2026-06-01 至 2026-07-30（数据实际覆盖
2026-06-01 至 2026-07-31）共 60 天的驾驶行为数据，预测每台车辆在未来 40 天内
是否发生事故（含因司机驾驶行为导致的未遂事故）。对每台车辆需输出：
\begin{enumerate}[label=(\arabic*)]
  \item \textbf{二分类预测结果}：未来 40 天内是否发生事故；
  \item \textbf{预测概率值}：0 至 1 之间的事故发生概率。
\end{enumerate}

可使用的驾驶行为特征维度包括：驾驶员状态（闭眼、困倦、打哈欠、打电话、看手机等）、
驾驶平稳性（超速、急加速、急刹车、风险入弯、长时间压线、前碰撞预警等）、超高危事件
（未遂事故、事故）、管理配合度（摄像头遮挡、角度异常）、环境道路（右转未停车、路口
超速、盲区预警）以及部分车辆的车载 IMU 数据。

官方评分标准为测试集上的 \textbf{AUC}（ROC 曲线下面积），即随机抽取一个正样本和一
个负样本时模型给正样本打分更高（含打平各计 0.5）的概率，按 AUC 对参赛团队排序。
因此本任务的全部建模工作以\textbf{提升折外 AUC 的排序能力}为核心目标，同时以
Recall@K、lift@K 等指标保证预警名单的可运营性。

\begin{hlbox}
\textbf{任务难点}（结合数据实测，供全组对齐认知）：
(1) 样本为车辆级且仅 500 台，正类稀缺（主口径下 126/500，正类占比 25.2\%），
高维稀疏行为特征下极易过拟合；
(2) 数据来自真实驾驶但经筛选、调整与模糊化处理，存在坐标占位符、传感器死区、
覆盖缺口等质量问题，需要系统化的清洗与可信度评估（见第 3 章）；
(3) 特征窗口（20 天）与预测窗口（40 天）的时序外推设定，要求训练与推理的
特征口径严格一致。
\end{hlbox}

\subsection{总体技术路线}

\begin{TODO}
此处放置全组技术路线总图（建议用 draw.io / PPT 绘制后导出 PNG/PDF 放入
\texttt{figures/} 目录，文件名建议 \texttt{pipeline.png}，正文用
%\includegraphics[width=\textwidth]{pipeline.png}
引用）。
\end{TODO}

路线概述（已按现有进展预写，建模线可修订）：原始四大类数据（车辆画像、风险事件、
轨迹、IMU）经\textbf{两阶段清洗}（第 3 章：一阶段单类型体检 + 二阶段跨类型互证，
产出 v0--v4 五版本数据）→ \textbf{特征工程}（第 4 章：事件计数三口径、轨迹暴露量、
IMU 覆盖与动作特征）→ \textbf{标签与划分}（第 5 章：20 天特征窗 + 40 天标签窗，
车辆级分层 5 折）→ \textbf{建模}（第 6 章：L2 逻辑回归 / 随机森林基线，
“单一模型”与“集成学习”双方案并行验证）→ \textbf{折外评估与预测生成}（第 7 章：
AUC 主指标 + 全组统一的实验准入规则）→ \textbf{提交}（第 8 章：CSV 结果文件 +
本说明文档 + 可复现代码）。

\begin{hlbox}
\textbf{分工与文档对应关系}：数据清洗与基础特征（第 2--4 章、附录 B）、
验证方案（第 5 章）、基线与建模（第 6--7 章）、总体协调与提交（第 1、8 章、附录 A/C）。
各章负责人见对应 \texttt{sections/*.tex} 文件头注释。
\end{hlbox}

\begin{TODO}
技术路线图（pipeline 总图）请建模线/组长补充；若全组对“单一模型 vs 集成学习”
双方案的最终取舍已有结论，请在 1.2 节末尾用一句话写明结论及第几章给出依据。
\end{TODO}
